\documentclass[11pt,letterpaper]{article}
\usepackage[utf8]{inputenc}
\usepackage{geometry}
\usepackage{setspace}
\usepackage{enumitem}

\geometry{left=1.0in,right=1.0in,top=0.75in,bottom=0.75in}
\setstretch{1.15}

\setlength{\parindent}{0pt}
\setlength{\parskip}{6pt plus 1pt minus 1pt}

\begin{document}

\thispagestyle{empty}

\begin{center}
    \large\textbf{IN THE UNITED STATES PATENT AND TRADEMARK OFFICE} \\
\end{center}

\vspace{0.2cm}
\noindent\begin{tabular}{p{2.3in}p{4.0in}}
\textbf{In re Patent Application of:} & Reinaldo Javier Aguiar Marcano \\
\textbf{Application Number:} & 19/560,038 \\
\textbf{Confirmation Number:} & 2359 \\
\textbf{Filing Date:} & March 7, 2026 \\
\textbf{Group Art Unit:} & Pre-Exam / TC 2100 \\
\textbf{Examiner:} & ADAM DEMARAY \\
\textbf{Title of Invention:} & OPTICAL COMPUTING ARCHITECTURE AND SYSTEM \\
\end{tabular}

\vspace{0.4cm}
\begin{center}
    \textbf{\large CONDITIONAL REQUEST FOR SUSPENSION OF ACTION \\
    UNDER 37 CFR \S{} 1.103(a)}
\end{center}

\vspace{0.2cm}
\noindent\textbf{Commissioner for Patents} \\
\noindent P.O. Box 1450 \\
\noindent Alexandria, VA 22313-1450

\vspace{0.2cm}
\noindent Sir:

Applicant Reinaldo Javier Aguiar Marcano, prosecuting the above-referenced nonprovisional patent application \textit{pro se} under Small Entity status, hereby formally requests a \textbf{six-month suspension of action} under 37 CFR \S{} 1.103(a) to take effect upon clearance of pre-examination formalities.

\section*{I. PROCEDURAL POSTURE}

Pursuant to the Office of Petitions Decision (OPPT) mailed October 1, 2026 (signed by Attorney Advisor Charles Steven Brantley), which noted that an applicant may seek relief under 37 CFR \S{} 1.103 after pre-examination processing, Applicant has concurrently filed a complete response curing all outstanding requirements from the August 5, 2026 Notice to File Missing Parts. Because the substitute specification, replacement claims and abstract on separate sheets, and statutory fees have been submitted, there will be no outstanding Office action once OPAP processes the reply, making a suspension under \S{} 1.103(a) procedurally proper upon transfer to Technology Center 2100.

\section*{II. GOOD AND SUFFICIENT CAUSE FOR SUSPENSION}

Applicant requests a six-month suspension of action on the merits for good and sufficient cause:

\begin{enumerate}[leftmargin=2em, itemsep=3pt, parsep=0pt]
    \item \textbf{Active Federal Civil Rights Litigation and Physical Confinement (SDTX Civil Action No. 4:26-CV-08794):} Applicant is currently prosecuting an active federal civil rights and emergency injunctive action in the United States District Court for the Southern District of Texas (Houston Division), docketed as \textbf{Civil Action No. 4:26-CV-08794}, \textit{Reinaldo J. Aguiar Marcano v. G. Troy Pickett, et al.} (assigned to \textbf{U.S. District Judge Kenneth M. Hoyt}). This federal action directly challenges and seeks to enjoin an unconstitutional state-court civil arrest warrant (Capias) issued in Fort Bend County, Texas (Cause No. 25-DCV-328122) without personal service or jurisdiction. Under the continuous threat of this void warrant, Applicant has been forced into an involuntary, \textit{de facto} physical confinement at his corporate office for over 145 consecutive days, completely stripping him of freedom of movement, preventing physical access to postal mail, restricting access to registered patent counsel, and severely incapacitating his ability to safely monitor and prosecute this patent application.
    \item \textbf{Pending Emergency Federal Injunctive Relief (TRO):} On October 7, 2026, the federal court officially entered Plaintiff's Complaint (Doc. 1) and Emergency \textit{Ex Parte} Application for a Temporary Restraining Order and Preliminary Injunction under Fed. R. Civ. P. 65 (Doc. 2). On October 8, 2026, Applicant submitted an Emergency Advisory and Request for Expedited Ruling to accelerate the federal injunction restraining county officials from executing the warrant.
    \item \textbf{Severe Financial Incapacitation:} The ongoing state-court asset freeze and physical confinement have subjected Applicant to a devastating financial burden. As a \textit{pro se} Small Entity inventor prosecuting a portfolio of over 100 related defensive technologies, Applicant faces aggregate progressive extension-of-time liabilities exceeding \$400,000. While Applicant has timely submitted the formalities cure and extension petition to protect this specific application today, sustaining continuous prosecution expenses across the portfolio while federal courts resolve his physical liberty is factually impossible.
\end{enumerate}

A six-month suspension of action will preserve the status quo, protect this critical technology from procedural forfeiture, and provide the federal judiciary the necessary time to adjudicate the emergency TRO and restore Applicant's physical capacity to prosecute.

\section*{III. FEE CALCULATION AND NOTICE OF USPTO PAYMENT GATEWAY OUTAGE}

The applicable Small Entity processing fee for a request for suspension under 37 CFR \S{} 1.17(i) (Fee Code 2830) is \textbf{\$60.00}.

\textbf{NOTICE OF PAYMENT PLATFORM OUTAGE (OCTOBER 10, 2026):} \\
Applicant notes that on today's filing date (Saturday, October 10, 2026), the USPTO electronic payment gateway (Patent Center / RAM / Financial Manager) is officially offline for scheduled system maintenance, as prominently announced via portal banner. Applicant submits this conditional request concurrently with the formalities cure today, and respectfully requests that the \$60.00 processing fee (Fee Code 2830) be assessed via formal Notice of Fee Deficiency under 37 CFR \S{} 1.53(f) and MPEP \S{} 506 in accordance with standard Office practice.

\vspace{0.3cm}
\begin{singlespace}
\noindent Respectfully submitted, \par\vspace{0.8cm}
\underline{\textbf{/Reinaldo Javier Aguiar Marcano/}} \\
\textbf{Reinaldo Javier Aguiar Marcano} \\
Pro Se Applicant \\
Mailing Address: 2302 Britton Ridge Drive, Katy, Texas 77494 \\
Telephone: (202) 599-3979 \\
Email: \texttt{rj.aguiar@gmail.com} (USPTO Account) \\
Email: \texttt{rj.aguiar.m@gmail.com} \\
Email: \texttt{reinaldo.aguiar.marcano@isbn9781971481296.com} \\
Date: October 10, 2026
\end{singlespace}

\end{document}
